% femtolisp 笔记
% license Usr
% type Note

\subsection{构建}
\begin{itemize}
\item 如果编译了 Julia，那么 \verb`src/flisp/flisp` 就是 REPL。
\item julia 中， \verb`julia/src` 或 \verb`julia/src/flisp` 有许多 scm 文件和 lsp 文件（都在 \verb`flsp` 中）。
\item \verb`julia-parser.scm` 是语法解析，\verb`ast.scm` 是语法树工具
\item \verb`aliases.scm` 中定义了许多 scheme 中的语法。 所以估计 scm 结尾的文件都需要先运行该文件才可以。
\end{itemize}


\subsection{类型}
\begin{itemize}
\item \href{https://code.google.com/archive/p/femtolisp/wikis/Manual.wiki}{语法手册（非官方）}
\item \verb`(typeof 1)` 是 \verb`fixnum`
\item \verb`(typeof 'abc)` 是 \verb`symbol`， 又如 \verb`'+`
\item \verb`'abc` 是 \verb`(quote abc)` 的语法糖
\item \verb`(typeof 1.23e4)` 是 \verb`double`
\item \verb`(typeof "abc")` 是 \verb`array byte`
\item \verb`(typeof #\a)` 是 \verb`wchar` 是字符
\item \verb`(typeof #t)` 是 \verb`boolean`， 还有 \verb`#f`
\item \verb`(typeof '())` 是 \verb`null`
\item \verb`(typeof (cons 1 2))` 是 \verb`pair`
\item \verb`(typeof (vector 1 2 3))` 是 \verb`vector`
\item \verb`(typeof (lambda (a b) (+ a b)))` 是 \verb`function`
\item \verb`(typeof aref)` 是 \verb`builtin` （内建函数）
\end{itemize}

\subsection{计算器}
\begin{itemize}
\item \verb`(+ 1 2 3)` 加法， 另有 \verb`-, *, /`
\item \verb`(define a 1)` 定义变量
\item \verb`(print 变量1 变量2 ...)` 打印变量值， \verb`princ` 也可以。 前者会给每个字符串添加双引号，后者不会，后者会把 \verb`\n` 打成换行符。 二者都会返回 \verb`#t`。 所以在 REPL 中后面会出现 \verb`#t` 并不是打印的内容而是返回值。
\item \verb`(define a 123)` 然后 \verb|`(1 ,a)| 会得到 \verb`(1 123)`。
\item \verb`(define a '(1 2 3))` 然后 \verb|`(0 ,@a 4)| 得到 \verb`(0 1 2 3 4)`。
\item \verb`(eq 1 1)` （\verb`eq` 是 \verb`eq?` 的别名）检查是否相等， 返回 \verb`#t`
\item \verb`(= 3 3.0)` 数值比较，即使是不同类型也能比。
\end{itemize}

\subsection{常用函数和 form}
\begin{itemize}
\item \verb`(cons 1 2)` 显示为 \verb`(1 . 2)`， \verb`(cons 1 (cons 2 (cons 3 4)))` 显示为 \verb`(1 2 3 . 4)`
\item \verb`(list 1 2 3 4)` 和 \verb`'(1 2 3 4)` 等效，显示为 \verb`(1 2 3 4)`
\item \verb`(length 列表或数组或字符串)` 返回长度（字符串返回字节数）， \verb`(length= 列表 长度)` 检查长度是否符合预期。
\item \verb`(aref (vector 10 11 12) 0)` 获取第 0 个元素 \verb`10`
\item \verb`(aset! 数组 索引 值)` 可以改变数组元素值。 如 \verb`(define x (vector 1 2 3 4))`， \verb`(aset! x 2 9)`
\item \verb`(car 变量)` 和 \verb`(cdr 变量)` 分别取第 1 和 2 个元素
\item \verb`(cadr 变量)` 等效于 \verb`(car (cdr 变量))`，以此类推。
\item \verb`(list 1 2 3)` 等效为 \verb`(cons 1 (cons 2 (cons 3 '())))`， 本质上是用二叉树表示单链表。 类型仍然是 \verb`pair`。
\item \verb`(append '(1 2) '(3 4))` 可以合并两个列表得到 \verb`(1 2 3 4)`
\item \verb`(append! my_list '(3 4))` 可以直接修改 \verb`my_list`。 例如 \verb`(define ll '(1 2 3))`， \verb`(append! ll '(4 5))`。 第二个参数不能是一个数字，一个数字也要写成 \verb`'(4)`
\item \verb`(set! 变量名 值)` 可以把 \verb`值` 赋值给变量名（可以不存在），并返回该值。 例如 \verb`(set! a (+ a 2))` 相当于 \verb`a += 2`。
\item \verb`(prog1 语句1 语句2 ...)` 依次执行每个语句并返回第一个语句的值。
\item \verb`(every 函数 列表)` 会检查函数作用在每个元素上是否都返回 \verb`#t`
\end{itemize}

\subsection{控制流}
\begin{itemize}
\item \verb`(if (> 2 1) 'big 'small)` 返回 \verb`'big`
\item 如果每个分支需要执行多个语句，用 \verb`(begin 语句1 语句2 ...)`， \verb`begin` 会返回最后一个语句。
\item 或者也可以用 \verb`(let () 语句1 语句2 ...)` 括号允许为空。
\begin{lstlisting}[language=none]
(if (> 3 2)
    (begin
      (princ "3 > 2\n")
      (princ "This is true branch\n"))
    (begin
      (princ "3 <= 2\n")
      (princ "This is false branch\n")))
\end{lstlisting}
\item 多个分支可以通过 if 的嵌套来表示
\begin{lstlisting}[language=none]
(define (grade-score s)
  (if (>= s 90)
      'A
      (if (>= s 80)
          'B
          (if (>= s 70)
              'C
              (if (>= s 60)
                  'D
                  'F)))))
\end{lstlisting}
又或者用 \verb`cond`
\begin{lstlisting}[language=none]
(define (grade-score s)
  (cond
    ((>= s 90) 'A)
    ((>= s 80) 'B)
    ((>= s 70) 'C)
    ((>= s 60) 'D)
    (else 'F)))
\end{lstlisting}
\item 另一个例子 \verb`(cond  ((= x 1) 'one)  ((= x 2) 'two) ((= x 3) 'three)  (else 'big))`
\end{itemize}

\subsubsection{循环}
\begin{itemize}
\item 循环只能用函数递归。例如计算 \verb`1+2+3+4`
\begin{lstlisting}[language=none]
(define (sum-up i sum)
  (if (> i 5)
      sum
      (sum-up (+ i 1) (+ sum i))))

(sum-up 1 0)
\end{lstlisting}
也可以用 \verb`let`
\begin{lstlisting}[language=none]
(let my-loop ((i 1) (sum 0))
  (if (> i 4)
      sum
      (my-loop (+ i 1) (+ sum i))))
\end{lstlisting}
还有一种等效写法（\verb`letrec` 允许定义的变量互相依赖，以及依赖自身）
\begin{lstlisting}[language=none]
(letrec ((loop-fun (lambda (i sum)
                       (if (> i 5)
                           sum
                           (loop-fun (+ i 1) (+ sum i))))))
  (loop-fun 1 0))
\end{lstlisting}

\item C++ 中的
\begin{lstlisting}[language=cpp]
int f() {
    int a = 1, b = 2, c = 3;
    for (int i = 0; i < 100; i++) {
        a = a + b;
        b = b + c;
        c = c + a;
    }
    return a + b + c;
}
\end{lstlisting}
翻译为
\begin{lstlisting}[language=none]
(define (f i a b c)
  (if (>= i 100)
      (+ a b c)
      (f (+ i 1)
         (+ a b)
         (+ b c)
         (+ c a))))
\end{lstlisting}

\item 把字符串变为字符列表（从后往前创建列表）
\begin{lstlisting}[language=none]
(define (string->list s)
  (let loop ((i (sizeof s)) (l '()))
    (if (= i 0)
	  l
	  (let ((newi (- i 1)))
	    (loop newi (cons (string.char s newi) l))))))
\end{lstlisting}
\end{itemize}

\subsubsection{异常}
\begin{itemize}
\item \verb`(trycatch (要执行的) 异常处理函数)`
\begin{lstlisting}[language=none]
(trycatch
  (begin
    (princ "Before error\n")
    (error "Something went wrong!")
    (princ "After error\n"))     ; never reached
  (lambda (e)
    (princ "Caught exception: " e "\n")))
\end{lstlisting}
\end{itemize}


\subsection{字符串}
\begin{itemize}
\item \verb`(string 变量)` 可以把其他类型转为字符串： \verb`string` 且支持多入参。 如 \verb`(string 1.2 ", " 3.4)` 得到 \verb`"1.2, 3.4"`
\item \verb`(length "你好吗")` 和 \verb`(sizeof "你好吗")` 返回字节数 9
\item \verb`(string.count "你好吗")` 返回字符数 3
\item \verb`(string "abc" "def" "ghi")` 拼接字符串
\item \verb`(string.join '("abc" "def") ", ")` 得到 \verb`"abc, def"`
\item \verb`(symbol "abc")` 把字符串转为符号。
\item \verb`(string.char 字符串 索引)` 获取某个字符。
\item \verb`(string.sub "abcde" 1 3)` 获取子字符串，得到 \verb`"bc"`
\item \verb`(string.inc 字符串 索引 字符数)` 向后跳 \verb`字符数` 个字符，返回新索引（编码 utf-8）。索引是字节为单位的。 \verb`索引` 可以不在字符起始位置。 \verb`string.dec` 同理
\item \verb`(string.find "abcdef" "cde")` 和 \verb`(string.find "abcdef" #\c)` 查找字符串， 都返回 2
\end{itemize}


\subsection{函数}
\begin{itemize}
\item 多次定义同名函数会把之前的定义覆盖。
\item \verb`(define (square n) (* n n))` 定义函数 \verb`(square 3)` 调用
\item \verb`(define (f1 a b) (+ a b))` 不完全是 \verb`(define f1 (lambda (a b) (+ a b)))` 的语法糖，函数对象会记录函数名为 \verb`f1`，无论赋值给谁。
\item \verb`(define f2 f1)` 可以赋值函数对象，但对象里面记录的函数名不会变。
\item \verb`(null? '())` 返回 \verb`#t`
\item \verb`(define (函数名 形参1 (形参2 默认值) ...) 语句1 语句2 ...)`。 如 \verb`(define (myfun1 x (y 3.14)) (+ x y))` 返回函数体中最后一个语句的值。
\item 返回 \verb`boolean` 的函数称为 predicate，习惯上以 \verb`?` 结尾
\item 内建 predicate（以下全返回 \verb`#t`）：\verb`(number? 42)`，\verb`(odd? 3)`， \verb`(string? "abc")`， \verb`(null? '())`，\verb`(pair? '(1 2))`，\verb`(even? 6)`，\verb`(eq? 1 1)`（判断是否同一对象），\verb`(eqv? 1 1)`（比较值和类型）。
\item \verb`(map 函数 列表)` 会把函数依次作用在列表的每个元素并返回一个新的列表。 例如 \verb`(map abs '(-1 -2 3 -4))` 得到 \verb`(1 2 3 4)`。 又如 \verb`(map (lambda (x) (* x x)) '(1 2 3))` 得到 \verb`(1 4 9)`。
\item \verb`*名字*` 只是个普通名字，习惯上用这种格式表示它是一个全局变量（声明在任何函数体之外）。
\item \verb`let` 结构
\begin{lstlisting}[language=none]
(let ((var1 val1)
      (var2 val2)
      ...)
  body-expr-1
  body-expr-2
  ...
  body-expr-n)
\end{lstlisting}
可以定义局部变量 \verb`var*`，并在后面的语句中使用，返回最后语句的结果。 返回后这些局部变量无定义。
\item \verb`let*` 可以让后定义的局部变量依赖于先声明的： \verb`(let* ((a 3) (b (+ a 1))) (* a b))` 返回 12。
\item 任意数量参数函数：\verb`(define (函数名 参数1 参数1 . 参数n) 函数体)`。例如 \verb`(define (show first . rest) (princ "First: " first "\n") (princ "Rest: " rest "\n"))`， 然后 \verb`(show 1 2 3 4)` 会显示 \verb`First: 1 Rest: (2 3 4)`。
\item 若要把列表展开作为调用函数的参数，如 \verb`(apply + '(1 2 3))` 得 6。
\item 列表是 immutable 的，若要改变，只能创建一个新的（也可以用 \verb`vector`）
\begin{lstlisting}[language=none]
(define (list-set lst n val)
  (if (null? lst)
      '()
      (if (= n 0)
          (cons val (cdr lst))
          (cons (car lst)
                (list-set (cdr lst) (- n 1) val)))))
\end{lstlisting}
使用： \verb`(list-set '(1 2 3 4) 2 99)` 得到 \verb`(1 2 99 4)`。
\end{itemize}

\subsection{文件读写}
\begin{itemize}
\item \verb`(file 文件名 :write :create)` 打开文件用于写入，若没有就创建。 返回一个 \verb`iostream` 类型，以下称为 \verb`流`。
\item \verb`:read` — open for reading.
\item \verb`:write` — open for writing.
\item \verb`:create` — create the file if it doesn’t exist.
\item \verb`:append` — open for appending (writing to end of file).
\item \verb`:truncate` — truncate (clear) existing file contents when opening for writing.
\item \verb`:text` — text mode (default).
\item \verb`:binary` — binary mode.
\item \verb`(io.write 流 "abc\n")` 写入并返回写入的字符数。
\item \verb`(io.close 流)` 关闭文件，若成功返回 \verb`#t`
\item \verb`(io.eof? 流)` 检查是否读到了文件末尾
\item \verb`(iostream? 变量)` 判断变量是否为流
\item \verb`(io.getc 流)` 读取一个字符，指针加一
\item \verb`(io.peekc 流)` 读取一个字符，指针不动
\item \verb`(io.putc 流)` 写入一个字符
\item \verb`(buffer)` 创建类似 \verb`std::sstream`，类型也是 \verb`iostream`，操作和文件流类似。
\item \verb`(io.seek 流 索引)` 移动指针到索引
\item \verb`(io.pos 流)` 返回当前索引
\end{itemize}


\subsection{宏}
\begin{itemize}
\item \verb`(define-macro (name arg1 arg2 ...) body...)` 定义一个宏。
\item 例如上面的 \verb`prog1` 宏， 若要重命名为 \verb`begin0`， 则用 \verb|(define-macro (begin0 first . rest) `(prog1 ,first ,@rest))|
\end{itemize}

